\documentclass[11pt]{article}

\usepackage[margin=1.05in]{geometry}
\usepackage{amsmath}
\usepackage{booktabs}
\usepackage{microtype}
\usepackage[hidelinks]{hyperref}
\setlength{\emergencystretch}{2em}

\title{Neither Synthetic Benchmarks nor Common-Sense Reasoning\\
Predicts Real-World Derivative-Free Optimizer Performance}
\author{Peter Cotton}
\date{\today}

\begin{document}
\maketitle

\begin{abstract}
Practitioners choose a derivative-free optimizer either by consulting a benchmark
leaderboard or by reasoning about the problem. We test both. Using a
memorisation-proof suite of real-world objectives, each relocated by a seeded
cube-to-cube diffeomorphism, we rank a panel of $22$ optimizers and ask how well that
ranking is predicted by the same panel's ranking on synthetic analytic benchmark
functions, and by a language model asked to reason, problem by problem, about which
optimizer should win. Both predictors are essentially uninformative: the rank
correlation with real-world performance is statistically indistinguishable from zero
in either case. A language-model selector that picks an optimizer per problem in fact
does worse than a constant policy of always using one robust direct-search method. The
two failures share a mechanism. Both the benchmark and the model over-trust
model-based trust-region methods (Powell's NEWUOA and BOBYQA), which dominate smooth
analytic functions yet rank near the bottom on the real suite. This is a draft, and
several confirmatory runs are still in progress (Section~\ref{sec:caveats}).
\end{abstract}

\section{The question}

Derivative-free optimization (DFO) has dozens of credible algorithms and no universal
winner. The two ways a practitioner actually picks one are to read a benchmark, such as
a global-optimization test-function leaderboard, or to reason about the problem (``it
is smooth and low-dimensional, so use a quadratic-model method''). We ask whether
either route predicts what actually works on real problems.

The difficulty in answering honestly is that public benchmarks are mostly synthetic
analytic functions such as sphere, Rosenbrock, and Rastrigin, and it is unknown whether
a leaderboard on those transfers to messy real objectives. We therefore build a
real-world instrument and treat the synthetic leaderboard and the model's reasoning as
two competing predictors of it.

\section{Instruments}

\paragraph{A test suite motivated by real problems.}
Standard global-optimization benchmarks are analytic test functions, chosen for known
optima and tunable difficulty rather than for resembling the problems practitioners
face. We instead use \texttt{example\_applications}, a collection of $52$ distinct
pure-Python objectives. Each is a reduced-order but faithful model of a genuine
engineering, scientific, or operational task: placing wind turbines to maximise output
under wake interaction, dispatching a grid battery against a price curve, throttling a
booster to a soft landing, arranging atoms to minimise the Lennard-Jones potential,
sizing a pressure vessel or cantilever, or tuning a PID controller, among
others.\footnote{Each problem has an interactive, visual demonstration at
\url{https://humpday.microprediction.org/applications/}.} Four of these expose a
faithful size knob (number of turbines, dispatch horizon, atom count, control segments)
that we scale to add $14$ higher-dimensional instances, so the full instance set spans
$2$ to $100$ dimensions, $66$ instances in all. These carry the structure, coupling,
ruggedness, and scale heterogeneity of real objectives, which is what analytic
benchmarks omit and what we are testing. All objectives live on $[0,1]^n$, and lower is
better.

\paragraph{Made memorisation-proof.}
Reusing a fixed set of named problems risks an optimizer, or its tuner, implicitly
learning where the optima are. We close this by wrapping each objective in a seeded
cube-to-cube diffeomorphism $\phi$, so the optimizer minimises $f(\phi^{-1}(x))$ on
$[0,1]^n$. This relocates the optimum to an instance-specific point while exactly
preserving the landscape's difficulty, conditioning, multimodality, and critical-point
structure; the value at each point is merely moved, not changed. Different seeds give
different but equally hard instances of the same problem, so nothing is won by
memorising a location, only by searching.

\paragraph{The panel and the empirical ranking.}
We score $22$ optimizers, namely the $21$ pure-Python \texttt{humpday} ports
(Nelder-Mead, Powell, three PRIMA trust-region methods, CMA-ES, Differential Evolution,
Particle Swarm, Pattern Search, Coordinate Descent, Simulated Annealing, and others)
together with production \texttt{pycma}, on each disguised instance at fixed evaluation
budgets. Per instance we rank the panel, and averaging gives a real-world leaderboard at
each budget.

\paragraph{Predictor 1: synthetic benchmarks.}
The same panel is scored on rotated synthetic test functions (sphere, ellipsoid, cigar,
Rastrigin, Rosenbrock) from the \texttt{nevergrad} library, at dimensions $5$, $20$, and
$40$, giving a synthetic leaderboard.

\paragraph{Predictor 2: model reasoning.}
For each real problem we give a language model only its natural-language description,
its dimension, and the evaluation budget, together with textbook one-line descriptions
of the candidate optimizers, and no benchmark results. It reasons about the likely
landscape and ranks the candidates. We repeat with shuffled presentation order and
aggregate the rankings with a Plackett-Luce (Bradley-Terry) model fit by Hunter's MM
algorithm; only the induced order is used, so cross-call score-scale arbitrariness never
enters.

\section{Result 1: synthetic benchmarks do not predict the real ranking}

Table~\ref{tab:e1} reports Kendall's $\tau$ and Spearman's $\rho$ between the synthetic
and real leaderboards, for the budgets completed so far. The correlation is weak and
not statistically significant.

\begin{table}[h]
\centering
\begin{tabular}{lccc}
\toprule
budget & Kendall $\tau$ & $p$ & Spearman $\rho$ \\
\midrule
60  & $+0.15$ & $0.32$ & $+0.21$ \\
120 & $+0.21$ & $0.17$ & $+0.27$ \\
\bottomrule
\end{tabular}
\caption{Synthetic-versus-real leaderboard correlation (synthetic $n=45$ instances,
real $n=90$). Both coefficients are indistinguishable from zero. These figures are
preliminary, from an earlier panel that included a Bayesian-optimization baseline since
removed for being intractable in high dimensions; they are being recomputed on the
current panel, and the budget-$240$ row is still in progress.}
\label{tab:e1}
\end{table}

The tails are where it bites. The synthetic top performers are the model-based methods
(NEWUOA, BOBYQA, UOBYQA, pycma, Powell), whereas the real-world top performers are
humble direct-search and population methods such as Coordinate Descent, Pattern Search,
Harmony Search, Ant Colony, and CMA-ES. NEWUOA ranks near the top of the panel on
synthetic functions, around second, but near the bottom on real problems, around
eleventh. The single best optimizer by synthetic benchmarks is one of the worst on real
problems.

\section{Result 2: model reasoning does no better, and is worse than a constant}

We elicited per-problem rankings for $18$ real problems at budget $120$. The predictive
correlation with the real ranking (Table~\ref{tab:llm}, left) is again about zero: the
model, at $\tau=+0.07$, is indistinguishable from random at $+0.03$ and only marginally
above synthetic at $-0.02$.

\begin{table}[h]
\centering
\begin{tabular}{lc@{\qquad}lc}
\toprule
\multicolumn{2}{c}{predictiveness ($\tau$ vs real)} & \multicolumn{2}{c}{selection: mean real rank of top pick} \\
predictor & $\tau$ & policy & mean rank \\
\midrule
model reasoning & $+0.07$ & oracle (best per problem) & $1.0$ \\
synthetic       & $-0.02$ & always Pattern Search & $3.78$ \\
random          & $+0.03$ & model's top pick per problem & $5.72$ \\
                &         & synthetic favourite (PRIMA) & $6.39$ \\
\bottomrule
\end{tabular}
\caption{Left: the model predicts the real ranking no better than chance. Right:
choosing the model's top pick per problem (lower is better, over a $10$-candidate field)
is worse than committing to one robust method.}
\label{tab:llm}
\end{table}

The selection test (Table~\ref{tab:llm}, right) is the sharper finding. Picking the
model's recommended optimizer per problem gives a mean real rank of $5.72$, worse than
always using a single robust method (Pattern Search, at $3.78$). The model's
problem-by-problem reasoning subtracts value relative to a constant policy.

The mechanism is explicit in the picks: $15$ of the $18$ top recommendations were PRIMA,
mostly BOBYQA. The model reliably reasons that a smooth-looking engineering objective
calls for a quadratic trust-region method, the very methods Result~1 shows crater on the
real suite, and it switches to CMA-ES or Differential Evolution only when the
description overtly signals ruggedness, as in the financial and game-playing problems.
The model reasons its way into the same trap the synthetic benchmark is built on.

\section{The shared mechanism}

Both failures have one cause. Smooth analytic functions reward an aggressively exploited
curvature model, so both the benchmark and textbook reasoning crown model-based
trust-region methods. Real objectives are rugged, coupled, noisy, or multi-scale, and at
low budget they reward methods that hedge rather than commit to a model. A separate probe
supports this directly: a covariance-damping intervention that reins in an over-trusted,
undersampled covariance helps on the structured real suite but is neutral to harmful on
smooth synthetic functions. It pays exactly when the model should not be trusted.

\section{Caveats and status}
\label{sec:caveats}

This is a living draft, and its numbers come with several qualifications. With a panel of
around two dozen optimizers and modest instance counts, none of the reported $\tau$
values is statistically significant; the honest claim is that synthetic benchmarks and
model reasoning carry no reliable signal about real-world performance, not that they are
perfectly anti-correlated. Several confirmatory runs are also still computing and will
harden or qualify these figures, including the budget-$240$ rank correlation, a
multi-budget held-out optimizer comparison, a runtime optimizer-crossover study, and a
train/test generalisation run.

Two further qualifications concern the baselines and the scope. The always-Pattern-Search
policy was the best fixed method on these problems, which gives it a mild post-hoc
advantage, though it is consistent with Pattern Search being a genuine top performer in
Result~1. And the model results rest on a single model, a single budget, and three
elicitation rounds, so they bound the effect rather than settle it.

\section{Implication}

If it holds, the practical message is uncomfortable and useful. One should not select a
DFO algorithm from a synthetic leaderboard, and should not trust a language model to
reason one out either, because both systematically over-recommend model-based methods
that fail on real problems. Until selection is grounded in real, or at least
memorisation-proof real-world-like, evaluation, committing to one robust and model-light
optimizer beats both. The broader caution is that a benchmark and a language model
trained on the same literature can share a blind spot, so agreement between them is not
evidence of correctness.

\end{document}
